\section{General spectral distributions}\label{sec:general}

The upper and lower bounds now determine the pressure for finite spectral
alphabets with rational proportions along the progressions used in the
cavity argument.  We first remove these arithmetic restrictions, and then
approximate an arbitrary compactly supported spectral law.  In this
section, subscripts in $R_\nu$ and $G_{t,\nu}$ specify the spectral law;
$\mathcal F(\nu)$ denotes the variational value in
\eqref{eq:variational} with that law.  The functions $S$ and $D_p$ do not
depend on the spectral law.

\subsection{Continuity of the spectral term}

The transform can reach the upper support edge, so its continuity must
include that branch of the definition.  The following statement needs
no regularity of the measure near the edge.

\begin{lemma}\label{gen:R-continuity}
Let $\nu_j$ and $\nu$ be probability measures supported in a common
interval $[-K,K]$.  Suppose that $\nu_j$ converges weakly to $\nu$ and
that
\[
 e_j:=\max\supp\nu_j\longrightarrow e:=\max\supp\nu.
\]
Define $b_{\nu_j}$ and $R_{\nu_j}$ by the edge convention of
Theorem~\ref{thm:main}, and likewise for $\nu$.  Then
\begin{equation}\label{gen:R-uniform}
 \sup_{0\le x\le1}|R_{\nu_j}(x)-R_\nu(x)|\longrightarrow0.
\end{equation}
Consequently,
\begin{equation}\label{gen:variational-continuity}
 |\mathcal F(\nu_j)-\mathcal F(\nu)|
 \le \frac12\sup_{0\le x\le1}|R_{\nu_j}(x)-R_\nu(x)|
 \longrightarrow0.
\end{equation}
\end{lemma}

\begin{proof}
For a measure $\tau$ with upper support edge $e_\tau$, put
\[
 H_\tau(z)=\int\frac{\tau(d\lambda)}{z-\lambda},
 \qquad z>e_\tau.
\]
This function is continuous and strictly decreasing.  Thus
$b_\tau(x)>e_\tau$ exactly when $x<H_\tau(e_\tau+)$, and in that
case $H_\tau(b_\tau(x))=x$.  In the remaining case
$b_\tau(x)=e_\tau$.  In both cases,
\begin{equation}\label{gen:b-bounds}
 e_\tau\le b_\tau(x)\le e_\tau+x^{-1},
 \qquad -K\le R_\tau(x)\le K.
\end{equation}
The first upper bound follows from
$H_\tau(z)\le(z-e_\tau)^{-1}$.  The lower bound on $R_\tau$
follows from $H_\tau(z)\ge(z+K)^{-1}$, using the edge limit of this
inequality when the transform is capped.

We first work away from zero.  Let $x_j\to x>0$ and pass to a
subsequence on which $b_{\nu_j}(x_j)\to b_*$.  The bounds
\eqref{gen:b-bounds} imply that such subsequences exist and that
$b_*\ge e$.  If $b_*>e$, the defining equations are uncapped for
large $j$.  Weak convergence, with their arguments uniformly separated
from the supports, gives
\[
 H_\nu(b_*)=x.
\]
Hence $b_*=b_\nu(x)$.  If $b_*=e$ but $b_\nu(x)>e$, choose
$c\in(e,b_\nu(x))$.  Then $H_\nu(c)>x$, so
$H_{\nu_j}(c)>x_j$ for large $j$.  Strict monotonicity, including the
definition of the capped branch, forces $b_{\nu_j}(x_j)>c$, a
contradiction.  These two cases prove
\[
 b_{\nu_j}(x_j)\longrightarrow b_\nu(x).
\]
The same argument with the measure fixed proves continuity of
$b_\nu$ on $(0,1]$.  A compactness argument now gives uniform
convergence of $R_{\nu_j}$ on every $[\varepsilon,1]$.

Near zero the convergence is uniform for a different reason.  If $K=0$
there is nothing to prove.  For $K>0$ and $x<(4K)^{-1}$, the transform
cannot be capped, since
$H_\tau(e_\tau+)\ge(2K)^{-1}$.  Write
$b_\tau(x)=x^{-1}+R_\tau(x)$ in its defining equation.  The bound
$|R_\tau|\le K$ permits the uniform expansion
\[
 1=\int\frac{\tau(d\lambda)}{1+x(R_\tau(x)-\lambda)}
  =1-x\left(R_\tau(x)-\int\lambda\,\tau(d\lambda)\right)
    +O(K^2x^2).
\]
Therefore
\begin{equation}\label{gen:R-at-zero}
 R_\tau(x)=\int\lambda\,\tau(d\lambda)+O(K^2x),
\end{equation}
with a constant independent of $\tau$.  Weak convergence on the common
bounded interval implies convergence of the means.  Combining
\eqref{gen:R-at-zero} with the convergence away from zero proves
\eqref{gen:R-uniform}.

Finally, $0\le D_p(r)\le1$ for every admissible $p$ and $r$, and hence
\[
 \sup_{p\in\cP}|G_{1,\nu_j}(p)-G_{1,\nu}(p)|
 \le\frac12\|R_{\nu_j}-R_\nu\|_\infty.
\]
Taking infima after adding the common functional $S$ proves
\eqref{gen:variational-continuity}.  These infima are finite: $S\ge0$
and $S(0)=0$ by Lemma~\ref{fld:domain}, whereas the transforms are
bounded by \eqref{gen:b-bounds}.
\end{proof}

\subsection{Spectral multiplicities and approximation}

For two diagonal spectra of the same size, using the same orthogonal
matrix gives the elementary comparison
\begin{equation}\label{gen:operator-comparison}
 |P_N(\Lambda,U)-P_N(\widetilde\Lambda,U)|
 \le\frac12\|\Lambda-\widetilde\Lambda\|_{\mathrm{op}}.
\end{equation}
It follows by comparing each spin energy and using $|\sigma|^2=N$.
The same argument gives monotonicity under coordinatewise increase of
the eigenvalues and the exact shift identity
\begin{equation}\label{gen:scalar-shift}
 P_N(\Lambda+c\Id,U)=P_N(\Lambda,U)+\frac c2.
\end{equation}
To change multiplicities, however, a small operator-norm error is not
available.  The next bound uses rotational symmetry inside each
eigenspace.

\begin{lemma}\label{gen:counts}
Fix distinct real numbers $\lambda_1,\ldots,\lambda_m$, and let
$L=\max_a\lambda_a-\min_a\lambda_a$.  Let $k_a$ and
$\widetilde k_a$ be positive integers whose sums are both $N$.
For diagonal matrices $\Lambda$ and $\widetilde\Lambda$ with these
multiplicities,
\begin{equation}\label{gen:count-comparison}
 \left|\E P_N(\Lambda,U)-\E P_N(\widetilde\Lambda,U)\right|
 \le\frac L2\max_a
 \frac{|k_a-\widetilde k_a|}{\min(k_a,\widetilde k_a)}.
\end{equation}
The same inequality holds for any fixed spin prior on the cube,
including a prior with deterministic external fields.
\end{lemma}

\begin{proof}
Arrange the two diagonal matrices so that, for each $a$,
$\min(k_a,\widetilde k_a)$ axes retain the value $\lambda_a$.
Haar invariance makes this rearrangement immaterial.  Set
$\Delta=\widetilde\Lambda-\Lambda$ and
\[
 f(s)=\E P_N(\Lambda+s\Delta,U),\qquad 0\le s\le1.
\]
Convexity of the log partition function gives
$f'(0)\le f(1)-f(0)\le f'(1)$.  At either endpoint,
\[
 f'(s)=\frac1{2N}\E\left\langle\sum_i\Delta_{ii}x_i^2
                  \right\rangle_s,
 \qquad x=U\sigma.
\]
Within a group of equal endpoint eigenvalues, the expected squared
projections are equal on all axes.  Indeed, rotating the corresponding
rows of $U$ preserves both its Haar law and the spin Hamiltonian; the
spin prior is unaffected.  If a group has $q_a$ axes and $c_a$ changed
axes, its changed axes therefore carry the fraction $c_a/q_a$ of the
group's expected squared spin norm.  Since $|\Delta_{ii}|\le L$ and
the total squared norm is $N$,
\[
 |f'(s)|\le\frac L2\max_a\frac{c_a}{q_a},
 \qquad s\in\{0,1\}.
\]
At the first endpoint $c_a=(k_a-\widetilde k_a)_+$, and at the second
$c_a=(\widetilde k_a-k_a)_+$.  The endpoint bounds prove
\eqref{gen:count-comparison}.
\end{proof}

\begin{proposition}\label{gen:approximation}
Let $\Lambda_N$ be deterministic diagonal matrices.  Suppose that their
empirical measures converge weakly to a compactly supported probability
measure $\mu$, and write
$a=\min\supp\mu$ and $b=\max\supp\mu$.  Assume
\[
 \liminf_{N\to\infty}\lambda_{\min}(\Lambda_N)\ge a,
 \qquad
 \limsup_{N\to\infty}\lambda_{\max}(\Lambda_N)\le b.
\]
Then $\E P_N\to\mathcal F(\mu)$.
\end{proposition}

\begin{proof}
We first finish the finite-alphabet case.  Suppose that
$\mu=\sum_{a=1}^m\rho_a\delta_{\lambda_a}$ with all $\rho_a>0$
rational.  Choose $n$ so that $s_a=n\rho_a$ are integers.  For each
residue $r\in\{0,\ldots,n-1\}$, choose nonnegative integers $c_a^r$
with sum $r$, and at dimension $N=kn+r$ use multiplicities
\[
 k_a(N)=ks_a+c_a^r.
\]
They increase by $s_a$ when $N$ increases by $n$, and are positive
for all sufficiently large $N$.  Proposition~\ref{up:bound} and
Proposition~\ref{cav:lower} give convergence to $\mathcal F(\mu)$ on
each progression.  There are finitely many residues, so these
canonically chosen multiplicities give convergence on the full
sequence of dimensions.

Now allow arbitrary positive proportions $\rho_a$ and arbitrary
multiplicities $k_a(N)$ with $k_a(N)/N\to\rho_a$.  Approximate
$\rho$ by positive rational probability vectors $\tau$, and let
$\widetilde k_a(N)/N\to\tau_a$ be the canonical counts just
constructed.  Lemma~\ref{gen:counts} gives
\begin{align*}
 &\limsup_{N\to\infty}
 \left|\E P_N(\Lambda_N,U)-
             \mathcal F\left(\sum_a\tau_a\delta_{\lambda_a}\right)
       \right|\\
 &\hspace{25mm}\le
 \frac L2\max_a\frac{|\rho_a-\tau_a|}{\min(\rho_a,\tau_a)}.
\end{align*}
The right side tends to zero as $\tau\to\rho$.
Lemma~\ref{gen:R-continuity} applies to these laws because their
support is the same fixed alphabet.  This proves the assertion for
every finite alphabet with positive limiting proportions.

We next construct finite alphabets for the general measure.  If $a=b$,
then $R_{\delta_a}=a$ and $\mathcal F(\delta_a)=a/2$ because
$S\ge0$ and $S(0)=0$; the assertion follows directly from
\eqref{gen:operator-comparison}.  Hence suppose $a<b$, and put
\[
 \eta_N=\max\{(a-\lambda_{\min}(\Lambda_N))_+,
                   (\lambda_{\max}(\Lambda_N)-b)_+\}.
\]
Then $\eta_N\to0$, and clipping the eigenvalues to $[a,b]$ changes
the pressure by at most $\eta_N/2$.  The clipped empirical measures
still converge weakly to $\mu$.

Fix $\delta>0$.  Partition $[a,b]$ into finitely many intervals of
length at most $\delta$, with all interior endpoints of zero
$\mu$-mass.  Use half-open intervals to assign endpoints uniquely,
including $a$ and $b$ in the first and last intervals.  For each
positive-mass interval $I_j$, choose
$r_j\in I_j\cap\supp\mu$, and put
\[
 \nu_\delta=\sum_{j:\,\mu(I_j)>0}\mu(I_j)\delta_{r_j}.
\]
The smallest and largest representatives satisfy
\[
 a\le r_{\min}\le a+\delta,
 \qquad b-\delta\le r_{\max}\le b,
\]
because every neighborhood of a support edge has positive mass.

In a positive-mass interval, replace each clipped eigenvalue by its
representative.  In all zero-mass intervals together, replace the
eigenvalues by $r_{\min}$ to form $\Lambda_N^-$, and by
$r_{\max}$ to form $\Lambda_N^+$.  Each zero-mass interval is a
continuity set for the clipped measures relative to $[a,b]$, so the
fraction of eigenvalues in their finite union tends to zero.  Both
modified spectra therefore have limiting law $\nu_\delta$, with
only positive limiting groups.  Moreover, coordinatewise,
\[
 \Lambda_N^--(\delta+\eta_N)\Id
 \le\Lambda_N\le
 \Lambda_N^++(\delta+\eta_N)\Id.
\]
Monotonicity and \eqref{gen:scalar-shift}, followed by the
finite-alphabet result, yield
\begin{equation}\label{gen:bin-squeeze}
 \mathcal F(\nu_\delta)-\frac\delta2
 \le\liminf_N\E P_N
 \le\limsup_N\E P_N
 \le\mathcal F(\nu_\delta)+\frac\delta2.
\end{equation}
Thus no estimate involving the fraction changed inside a vanishing
group is needed.

As $\delta\downarrow0$, the laws $\nu_\delta$ converge weakly to
$\mu$ and their upper support edges converge to $b$.
Lemma~\ref{gen:R-continuity} gives
$\mathcal F(\nu_\delta)\to\mathcal F(\mu)$, and
\eqref{gen:bin-squeeze} completes the proof.  The thermodynamic limit
was taken for each fixed finite alphabet before its refinement.
No derivative bound on the finite-alphabet transforms is required
uniformly over those refinements.
\end{proof}

\begin{proof}[Completion of the proof of Theorem~\ref{thm:main}]
Proposition~\ref{gen:approximation} supplies convergence of the
expected pressures.  The spectral hypotheses also imply
$\|\Lambda_N\|_{\mathrm{op}}\le K$ for all sufficiently large $N$,
for some finite $K$.  Comparing the quadratic forms for two orthogonal
matrices shows that $P_N(\Lambda_N,U)$ is $K$-Lipschitz in the
Frobenius distance, independently of the number of distinct
eigenvalues.  The fourth-moment argument in
Lemma~\ref{rot:concentration} therefore gives
\[
 \E|P_N-\E P_N|^4\le \frac{C_K}{N^2}.
\]
For every $\varepsilon>0$, Markov's inequality makes the probabilities
of $|P_N-\E P_N|>\varepsilon$ summable.  Borel--Cantelli, first for
positive rational $\varepsilon$, proves $P_N-\E P_N\to0$ almost
surely.  This argument needs no independence between the orthogonal
matrices at different dimensions.
\end{proof}

The pressure limit now also gives consequences for random spectra and
maximum energies. We derive these first, then apply them to
Gaussian-pattern interactions.

\subsection{Random spectra and zero temperature}

\begin{proposition}\label{gen:random}
Let $J_N$ be random real symmetric matrices whose conditional law,
given their eigenvalues, is the orthogonal orbit measure.  Let $\mu$
be a fixed compactly supported probability measure.  If the empirical
spectral measures converge weakly to $\mu$ in probability and the two
edge excesses in Proposition~\ref{gen:approximation} tend to zero in
probability, then $P_N\to\mathcal F(\mu)$ in probability.  If these
spectral hypotheses hold almost surely, the pressure convergence
holds almost surely.  In either case, if the random variables $P_N$
are uniformly integrable, then $\E P_N\to\mathcal F(\mu)$.
A sufficient condition for this last hypothesis is uniform
integrability of $\|J_N\|_{\mathrm{op}}$, in particular a uniform
deterministic bound on these norms.
\end{proposition}

\begin{proof}
Let $f_N(\Lambda)$ be the Haar expectation of the pressure with
deterministic eigenvalues $\Lambda$.  The integral against the
conditional orbit law is therefore the finite quantity
$f_N(\Lambda_N)$.  This is used as a conditional orbit integral,
without any assumption of global integrability.  Equivalently, all
conditional calculations below may first be restricted to
$\{\|\Lambda_N\|_{\mathrm{op}}\le K\}$, on which $|P_N|\le K/2$.
Under the almost-sure spectral hypotheses,
Proposition~\ref{gen:approximation} applies to each spectral sequence
in the event of convergence and yields
$f_N(\Lambda_N)\to\mathcal F(\mu)$ almost surely.  Under the
in-probability hypotheses, every subsequence has a further subsequence
on which all the spectral hypotheses hold almost surely.  Here weak
convergence can be expressed by a metric on probability measures, so
only finitely many convergence-in-probability assertions are involved.
The same deterministic result on this further subsequence proves
$f_N(\Lambda_N)\to\mathcal F(\mu)$ in probability.

For each fixed $K$, conditional concentration gives
\begin{equation}\label{gen:localized-concentration}
 \Prob\bigl(|P_N-f_N(\Lambda_N)|>\varepsilon,
                   \|\Lambda_N\|_{\mathrm{op}}\le K\bigr)
 \le\frac{C_K}{\varepsilon^4N^2}.
\end{equation}
The spectral hypotheses in probability imply that the norm bound holds
with probability tending to one for a sufficiently large fixed $K$.
This proves convergence of the fluctuations in probability.  Under
the almost-sure hypotheses the spectra are eventually bounded almost
surely.  Apply Borel--Cantelli to
\eqref{gen:localized-concentration} for integer $K$ and positive
rational $\varepsilon$ to obtain almost-sure convergence of the
fluctuations as well.

Uniform integrability and convergence in probability imply convergence
of expectations.  Finally,
$|P_N|\le\|J_N\|_{\mathrm{op}}/2$, which proves the sufficient
condition.
\end{proof}

\begin{corollary}\label{gen:ground-state}
Let $J_N=U_N^{\mathsf T}\Lambda_NU_N$ satisfy the deterministic
spectral hypotheses of Theorem~\ref{thm:main}, with limiting law
$\mu$.  Define its maximum energy per spin by
\[
 M_N=\frac1{2N}\max_{\sigma\in\{-1,1\}^N}
                         \sigma^{\mathsf T}J_N\sigma.
\]
For $\beta>0$, let $\mu_\beta$ be the image of $\mu$ under
$\lambda\mapsto\beta\lambda$.  Then, almost surely and in
expectation,
\begin{equation}\label{gen:ground-state-formula}
 M_N\longrightarrow
 \lim_{\beta\to\infty}\frac{\mathcal F(\mu_\beta)}\beta.
\end{equation}
In particular, the limit on the right exists.
\end{corollary}

\begin{proof}
Write $P_N(\beta)$ for the pressure with spectrum
$\beta\Lambda_N$.  The uniform probability on the cube gives, for
every $N$ and $\beta>0$,
\begin{equation}\label{gen:maximum-squeeze}
 \frac{P_N(\beta)}\beta
 \le M_N\le
 \frac{P_N(\beta)}\beta+\frac{\log2}\beta.
\end{equation}
Theorem~\ref{thm:main} applies at each positive integer $\beta$.
Intersecting these countably many probability-one events, the
difference between the limit superior and limit inferior of $M_N$ is
at most $\log2/\beta$ for every such integer.  Thus $M_N$ converges
almost surely.  Moreover, its limit belongs to each deterministic
interval
\[
 \left[\frac{\mathcal F(\mu_\beta)}\beta,
       \frac{\mathcal F(\mu_\beta)}\beta+\frac{\log2}\beta\right],
 \qquad \beta\in\N,
\]
so it is a deterministic number $M$.  The eventual deterministic bound
on $\|J_N\|_{\mathrm{op}}$ also gives convergence in expectation.
For any real $\beta>0$, take expectations and then $N\to\infty$ in
\eqref{gen:maximum-squeeze}.  It follows that
\[
 0\le M-\frac{\mathcal F(\mu_\beta)}\beta
 \le\frac{\log2}\beta,
\]
which proves \eqref{gen:ground-state-formula} without a uniform
positive-temperature estimate.
\end{proof}

\subsection{Gaussian-pattern Wishart interactions}

Fix $\alpha>0$, put $m_N=\lfloor\alpha N\rfloor$, and let $G_N$ be an
$m_N$-by-$N$ matrix of independent standard real Gaussian variables.
For $c\in\R$, define
\begin{equation}\label{gen:wishart-pressure}
 p_N(c)=\frac1N\log\left[
 2^{-N}\sum_{\sigma\in\{-1,1\}^N}
 \exp\left\{\frac{c}{2N}\|G_N\sigma\|_2^2\right\}\right].
\end{equation}
This is the Gaussian-pattern Hopfield pressure for
$J_N=cG_N^{\mathsf T}G_N/N$, with uniform spin probability and the full
quadratic form, including its diagonal self-interactions.

Let $a_\alpha=(1-\sqrt\alpha)^2$, $b_\alpha=(1+\sqrt\alpha)^2$, and set
\begin{equation}\label{gen:wishart-law}
 \begin{gathered}
 \mu_\alpha=(1-\alpha)_+\delta_0+
 \frac{\sqrt{(b_\alpha-x)(x-a_\alpha)}}{2\pi x}
 \ind_{[a_\alpha,b_\alpha]}(x)\,dx,\\
 \mu_{\alpha,c}=(x\mapsto cx)_\#\mu_\alpha.
 \end{gathered}
\end{equation}
The displayed density is understood almost everywhere.

\begin{corollary}[Gaussian patterns]\label{gen:wishart}
For every fixed $c\in\R$,
\begin{equation}\label{gen:wishart-limit}
 p_N(c)\longrightarrow\mathcal F(\mu_{\alpha,c})
\end{equation}
in probability and in $L^1$, and hence in expectation. The subsequent
zero-temperature limits
\begin{equation}\label{gen:wishart-zero-temperature}
 e_\pm(\alpha):=
 \lim_{\beta\to\infty}\frac{\mathcal F(\mu_{\alpha,\pm\beta})}{\beta}
\end{equation}
exist, with $e_+(\alpha)\ge0$ and $e_-(\alpha)\le0$. If
$g_{N,1},\ldots,g_{N,N}\in\R^{m_N}$ are the columns of $G_N$, then
\begin{align}
 \frac{\|G_N\|_{\infty\to2}}N
 &=\frac{\|G_N^{\mathsf T}\|_{2\to1}}N
 \longrightarrow\sqrt{2e_+(\alpha)}, \label{gen:gaussian-max-norm}\\
 \frac1N\min_{\sigma\in\{-1,1\}^N}
 \left\|\sum_{i=1}^N\sigma_i g_{N,i}\right\|_2
 &\longrightarrow\sqrt{-2e_-(\alpha)}. \label{gen:gaussian-min-sum}
\end{align}
Both limits hold in probability and in $L^1$. The operator norms are
from $\ell_\infty^N$ to $\ell_2^{m_N}$ and from $\ell_2^{m_N}$ to
$\ell_1^N$, respectively.
\end{corollary}

\begin{proof}
Put $W_N=G_N^{\mathsf T}G_N/N$ and $\alpha_N=m_N/N$. For every
$Q\in O(N)$, the matrices $G_NQ$ and $G_N$ have the same law.
Haar averaging this identity against any bounded function of the
ordered eigenvalues shows that the conditional law of $cW_N$ is its
orthogonal orbit law, including when the zero eigenvalue is repeated.

The Marchenko--Pastur theorem
\cite[Theorem~1 and the examples on pp.~459 and 461]{MarchenkoPastur1967}
gives weak convergence in probability of the empirical law of $W_N$
to $\mu_\alpha$. To check the edges, let $s_{\min}^+(G_N)$ be the
smallest positive singular value. For $m_N\ge1$ and $t>0$, the Gaussian
singular-value bounds \cite[Theorem~II.13]{DavidsonSzarek2001}, after
scaling and transposing when necessary, give
\begin{align*}
 \Prob\left(\frac{s_{\max}(G_N)}{\sqrt N}>
        1+\sqrt{\alpha_N}+t\right)&\le e^{-Nt^2/2},\\
 \Prob\left(\frac{s_{\min}^+(G_N)}{\sqrt N}<
        |1-\sqrt{\alpha_N}|-t\right)&\le e^{-Nt^2/2}.
\end{align*}
Together with positive limiting mass near each continuous spectral edge,
these bounds give $\lambda_{\max}(W_N)\to b_\alpha$ in probability.
The lower edge is $0$ for $\alpha\le1$ and $a_\alpha$ for $\alpha>1$.
For $\alpha<1$, $W_N$ has $N-m_N$ zero eigenvalues almost surely for
all large $N$; at $\alpha=1$, nonnegativity and limiting mass near zero
give the lower-edge limit; for $\alpha>1$, the second bound and limiting
mass near $a_\alpha$ give that limit. Thus both edge excesses vanish.
Scaling by a fixed $c$ preserves this conclusion, reversing the edges
when $c<0$; the case $c=0$ is immediate.

Since $\alpha_N\le\alpha$, integrating the first singular-value bound
gives
\[
 \sup_N\E\|W_N\|_{\mathrm{op}}^2
 =\sup_N\E\left(\frac{s_{\max}(G_N)}{\sqrt N}\right)^4<\infty.
\]
The finitely many possible cases $m_N=0$ are zero matrices. Hence
$\|cW_N\|_{\mathrm{op}}$ and $p_N(c)$ are uniformly integrable for
each fixed $c$. Proposition~\ref{gen:random} gives convergence in
probability, and uniform integrability upgrades it to $L^1$.

For $\varepsilon\in\{+1,-1\}$, put
\[
 H_N^\varepsilon=\frac1{2N}
 \max_{\sigma\in\{-1,1\}^N}\sigma^{\mathsf T}(\varepsilon W_N)\sigma,
 \qquad
 A_\beta^\varepsilon=
 \frac{\mathcal F(\mu_{\alpha,\varepsilon\beta})}{\beta}.
\]
The same maximum squeeze as in \eqref{gen:maximum-squeeze} holds
pointwise:
\[
 \frac{p_N(\varepsilon\beta)}{\beta}
 \le H_N^\varepsilon\le
 \frac{p_N(\varepsilon\beta)}{\beta}+\frac{\log2}{\beta}.
\]
Also $|H_N^\varepsilon|\le\|W_N\|_{\mathrm{op}}/2$, so these energies
are uniformly integrable. Taking expectations and then $N\to\infty$
at each fixed $\beta>0$ gives
\[
 A_\beta^\varepsilon
 \le\liminf_N\E H_N^\varepsilon
 \le\limsup_N\E H_N^\varepsilon
 \le A_\beta^\varepsilon+\frac{\log2}{\beta}.
\]
The expectations are bounded, and the gap between their limit superior
and limit inferior is at most $\log2/\beta$ for every $\beta>0$.
They therefore have a common limit $e_\varepsilon(\alpha)$, and the
same bounds give
\[
 0\le e_\varepsilon(\alpha)-A_\beta^\varepsilon
 \le\frac{\log2}{\beta},
\]
which proves \eqref{gen:wishart-zero-temperature}. Given any error
tolerance, choose one sufficiently large fixed $\beta$ and use
\eqref{gen:wishart-limit} in the pointwise squeeze. This proves
$H_N^\varepsilon\to e_\varepsilon(\alpha)$ in probability; uniform
integrability gives convergence in $L^1$.

Dualizing the Euclidean norm and maximizing a linear functional over
the cube gives
\[
 \|G_N\|_{\infty\to2}
 =\|G_N^{\mathsf T}\|_{2\to1}
 =\max_{\sigma\in\{-1,1\}^N}\|G_N\sigma\|_2.
\]
Consequently,
\[
 H_N^+=\frac{\|G_N\|_{\infty\to2}^2}{2N^2},
 \qquad
 H_N^-=-\frac1{2N^2}\min_{\sigma\in\{-1,1\}^N}\|G_N\sigma\|_2^2.
\]
The signs of $e_+$ and $e_-$ follow from these identities. Since
$G_N\sigma=\sum_i\sigma_i g_{N,i}$, the $L^1$ convergence of the
energies and the inequality
$|\sqrt{x}-\sqrt{y}|\le\sqrt{|x-y|}$ for $x,y\ge0$ give
\eqref{gen:gaussian-max-norm}--\eqref{gen:gaussian-min-sum}.
\end{proof}

For comparison, the finite partition expression in the companion
Ising-perceptron manuscript with $f_c(z)=cz^2/2$ is exactly
\eqref{gen:wishart-pressure}. For a standard Gaussian $Z$,
$\E|f_c(Z)|=|c|/2$, while
$\E e^{f_c(Z)}=(1-c)^{-1/2}$ for $c<1$ and is infinite for $c\ge1$.
Its integrable-activation theorem
\cite[Theorem~8.3]{OpenAIIsingPerceptron2026} therefore applies precisely
when $c<1$ in this quadratic family. Uniqueness of the expected-pressure
limit gives the equality of variational values
\[
 \mathcal F(\mu_{\alpha,c})=\mathcal P_{f_c}(\alpha),
 \qquad c<1.
\]
This compares values, not the functionals pointwise; the companion
functional is not invoked for $c\ge1$.

\subsection{Why weak convergence on an unbounded support is insufficient}

The bounded-edge hypothesis cannot be replaced by the assertion that
all eigenvalues lie in the support of a possibly unbounded limiting
law.  We give an explicit sequence for which that weaker condition
holds but the pressure has no limit.

Let $\mu$ have density $e^\lambda\ind_{\{\lambda\le0\}}$.
For every $N$, its midpoint quantiles are
\[
 q_{N,j}=\log\frac{j-1/2}{N},\qquad 1\le j\le N.
\]
Their empirical measures converge weakly to $\mu$.  At even $N$ use
these eigenvalues.  At odd $N$ replace $q_{N,1}$ by $-e^{N^2}$ and
leave the other eigenvalues unchanged.  Changing one eigenvalue does
not affect weak convergence, and every eigenvalue still belongs to
$\supp\mu=(-\infty,0]$.

At even sizes, convexity of $-\log u$ gives the midpoint bound
\[
 \frac1N\sum_{j=1}^N(-q_{N,j})
 \le\int_0^1-\log u\,du=1.
\]
Jensen's inequality for the uniform spin measure, whose mean quadratic
energy is $\Tr\Lambda_N/2$, therefore yields
\begin{equation}\label{gen:even-counterexample}
 -\frac12\le\frac{\Tr\Lambda_N}{2N}\le P_N\le0.
\end{equation}

At odd sizes, let $v_N$ be the unit eigenvector corresponding to
$-e^{N^2}$.  For each fixed spin vector $\sigma$, Haar invariance
implies that $v_N\cdot\sigma$ has the distribution of $\sqrt N$
times the first coordinate of a uniform point on the unit sphere in
$\R^N$.  For $N\ge3$ its density on a fixed neighborhood of zero is
bounded by an absolute constant.  Indeed, that density is
\[
 \frac{\Gamma(N/2)}{\sqrt{\pi N}\,\Gamma((N-1)/2)}
 \left(1-\frac{t^2}{N}\right)^{(N-3)/2},
 \qquad |t|<\sqrt N,
\]
whose prefactor is bounded.  A union bound over the cube gives
\[
 \Prob\left(\min_{\sigma\in\{-1,1\}^N}
                  |v_N\cdot\sigma|<e^{-2N}\right)
 \le C2^Ne^{-2N}.
\]
This is summable.  Almost surely, for all sufficiently large odd $N$,
every spin vector therefore has energy at most
$-\tfrac12e^{N^2-4N}$, since all remaining eigenvalues are
nonpositive.  Consequently,
\begin{equation}\label{gen:odd-counterexample}
 P_N\le-\frac{e^{N^2-4N}}{2N}\longrightarrow-\infty
 \qquad\text{along odd }N.
\end{equation}
The bounds \eqref{gen:even-counterexample} and
\eqref{gen:odd-counterexample} prove the claimed failure of
convergence.  No independence between dimensions is used here either.
